\chapter{Conclusions}
\label{ch:conclusions}

\begin{guidance}[What this chapter does]
Answers the aim and each objective from \cref{sec:objectives}, drawing only
on evidence already presented -- no new material. The Handbook suggests
presenting conclusions in an itemised form. A pattern that works well is to
repeat each objective and summarise below it what was achieved and where the
evidence is (adapted from a PhD thesis by \textcite{bausch2014}). Typical length:
400--600 words.
\end{guidance}

\section{Resolution of the aim and objectives}
\label{sec:resolution}

The aim of the project was to \ldots\ This aim was \ldots

\begin{description}[style=nextline, leftmargin=1.5em]
  \item[\Cref{obj:review}: To review existing solutions and the relevant literature \ldots]
    Met. \Cref{ch:literature} reviewed \ldots\ and identified \ldots
  \item[\Cref{obj:spec}: To produce a product design specification \ldots]
    Met. The specification in \cref{ch:specification} defined five
    requirements, which were refined after \ldots
  \item[\Cref{obj:design}: To generate and evaluate candidate solutions \ldots]
    Met. Three concepts were evaluated with a weighted decision matrix
    (\cref{tab:decision}) \ldots
  \item[\Cref{obj:build}: To implement and test a prototype \ldots]
    Partly met. The prototype passed two of four tests (\cref{tab:tests});
    \cref{req:cost} was not met because \ldots
  \item[\Cref{obj:evaluate}: To evaluate the prototype and the process \ldots]
    Met. \Cref{ch:evaluation} \ldots
\end{description}

\section{Overall conclusion}
\label{sec:overall-conclusion}
